\centerline{\bf \Huge Солянка II}
\centerline{\bf \Huge }

\begin{flushright}
        \scriptsize{}
\end{flushright}

\problem{1}
Существует ли многочлен $P(x)$ такой, что $P(20)=20, P(26)=26$ и $$\forall \ n \in \mathbb{N}\backslash\{20,26\}:  P(n)\not\in \mathbb{Q}$$

%бибиков

\problem{2}
Дано нечетное простое число $p$. Чему может быть равен НОД чисел $a+b$ и $\frac{a^p+b^p}{a+b}$, если $a$ и $b$ - взаимно простые целые числа?

%бибиков

\textit{Подсказка:} $4G, H+$

\problem{3}
Пусть $A B C$ треугольник, в котором $\angle A C B=90^{\circ}$ и $\angle A>\angle B$. Касательная в точке $C$ к описанной окружности треугольника $A B C$, пересекает прямую $A B$ в точке $D$. Точка $E$ середина отрезка $C D$, а точка $F$ лежит на прямой $E B$ так, что $A F \| C D$. Докажите, что прямые $A B$ и $C F$ перпендикулярны.

%колм

\textit{Комментарий.} Подумайте над тем, какое есть обобщение у задачи.

\problem{4}
В клетках квадрата $10 \times 10$ расставлены натуральные числа от 1 до 100 , каждое по одному разу. Эльзята рассмотрела всевозможные квадраты $2\times 2$ со сторонами, идущими по линиям сетки, и в каждом отметила наибольшее число (при этом одно число могло быть отмечено несколько раз). Могло ли так оказаться, что все двузначные числа были отмечены?

%https://math.mosolymp.ru/upload/files/2025/khamovniki/9/2025-05-15_Otborochnaya_olimpiada,_resheniya.pdf